\documentclass[conference,nofonttune]{IEEEtran}
\IEEEoverridecommandlockouts

% Keep the IEEEtran layout but use Computer Modern instead of its Times/
% Helvetica/Courier defaults (IEEEtran.cls forces ptm/phv/pcr). lmodern is
% Computer Modern shipped as complete Type 1 fonts (with T1 encoding it also
% covers glyphs like the section sign, avoiding Type 3 bitmap fallbacks); it
% overrides the class font defaults. Math already uses Computer Modern. The
% nofonttune class option above disables IEEEtran's Times-specific interword
% tuning.
\usepackage[T1]{fontenc}
\usepackage{lmodern}

\usepackage{cite}
\usepackage{amsmath,amssymb,amsfonts}
\usepackage{booktabs}
\usepackage{graphicx}
\usepackage{xcolor}
\usepackage{listings}
\usepackage[hidelinks]{hyperref}
\usepackage{textcomp}

\lstset{
  basicstyle=\ttfamily\footnotesize,
  breaklines=true,
  columns=fullflexible,
  keepspaces=true,
  frame=single,
  framesep=3pt,
  xleftmargin=3pt,
  showstringspaces=false,
}

\newcommand{\code}[1]{\texttt{#1}}
\newcommand{\reg}[1]{\texttt{#1}}

\begin{document}

\title{Booting an Operating System on a RISC-V Soft Core:\\
A Development History from RV32/SERV to an\\
Interrupt-Driven Zircon on NOEL-V}

\author{\IEEEauthorblockN{Filip Filmar}
\IEEEauthorblockA{\textit{Independent} \\
filmil@gmail.com}}

\maketitle

\begin{abstract}
We chronicle the multi-month bring-up of a full operating-system stack on a
RISC-V soft core on a single low-cost FPGA: from a bit-serial RV32 bootstrap
loader, through core, UART, and DRAM bring-up on a 64-bit GRLIB NOEL-V core, to
M-mode firmware (OpenSBI), a supervisor-mode boot ladder, and finally a
fully interrupt-driven console for Google's Zircon microkernel. The value of the
account is less any single fix than the \emph{pattern}: almost every failure
traced not to the component under suspicion but to a non-standard choice in our
own environment---an unconfigured scan input, a 32-bit core running 64-bit code,
a byte-write strobe hardwired on, a device tree that under- or over-declared RAM,
an interrupt-controller that was correct all along, and two RISC-V extensions
(\emph{Smrnmi}, \emph{Sstc}) that were present in silicon but never enabled or
advertised. We describe a three-vantage debugging methodology---isolated IP
simulation, non-perturbing bare-metal probes, and a JTAG debug module that reads
CSRs, arbitrary memory, and the kernel's own in-RAM log---that repeatedly
separated real hardware faults from measurement artifacts and firmware
omissions, and that made a low-cost, GPL-only toolchain sufficient to debug a
proprietary core booting a production microkernel.
\end{abstract}

\begin{IEEEkeywords}
RISC-V, NOEL-V, SERV, Zircon, Fuchsia, OpenSBI, PLIC, Smrnmi, Sstc, DDR3,
FPGA, JTAG, bring-up, debugging.
\end{IEEEkeywords}

\section{Introduction}

Soft-core RISC-V on an FPGA is uniquely inspectable: RTL, firmware, and kernel
are all editable, and a debug module can read any register or memory word on the
live system. That same openness means a single observable symptom---a core that
fetches garbage, a store that ``never reaches the bus,'' a console that goes
quiet---can originate in any layer, and the layers interact in ways that defeat
single-layer reasoning. This paper is a development history of exactly that
situation, assembled from the working notes of the bring-up.

The target is Zircon, the microkernel under Google's Fuchsia, running on a
NOEL-V core (from the GRLIB IP library) synthesized to a Xilinx Artix-7
(xc7a200t) at 40\,MHz. Reaching a booting kernel required climbing a long ladder,
and the recurring lesson---which the developer's notes phrase as ``why doesn't
everyone hit this?''---is that \emph{every} obstacle was a consequence of a
non-standard decision in the build, not a defect in the widely used component we
first blamed. We present the arc in the order it was climbed:

\begin{itemize}
  \item \S\ref{sec:bootstrap}: the RV32/SERV bootstrap that loads and launches the 64-bit core;
  \item \S\ref{sec:core}: core bring-up---an X-propagation that poisoned the whole pipeline in simulation, a core accidentally built 32-bit, and a UART that worked in simulation but not on silicon;
  \item \S\ref{sec:mem}: the DRAM subsystem---a byte-write strobe hardwired on, and a 128\,MB aliasing bug that masqueraded as a memory-size limit;
  \item \S\ref{sec:opensbi}: M-mode firmware---OpenSBI wedged by heap corruption and by load-reserved/store-conditional that this core does not hold on uncached DRAM;
  \item \S\ref{sec:ladder}: the supervisor-mode boot ladder to a full Zircon boot, including a shim that trapped inside a bounds-checked C++ iterator and a data image placed on top of the boot stub;
  \item \S\ref{sec:console}: the interrupt-driven console, whose five-layer root cause is the technical climax; and
  \item \S\ref{sec:method}: the debugging methodology that made all of it tractable.
\end{itemize}

A timeline of the bring-up, drawn from the project's commit history, appears in
Appendix~\ref{app:timeline}.

\section{The Bootstrap Architecture}
\label{sec:bootstrap}

The board carries no boot ROM and no bootloader flash. Bring-up therefore began
not with the NOEL-V but with \textbf{SERV}, a bit-serial RV32 core small enough
to bootstrap from block RAM. Early boards (\code{rv32\_uart},
\code{rv32\_ddr3}, \code{rv32\_ddr3exec}) established, incrementally, that the
RV32 core could drive the UART, reach DDR3, and \emph{execute} from DDR3.

That RV32 core then became the permanent first stage. In the production design
the SERV core boots, runs an Intel-HEX \emph{uploader} over the serial port, and
streams an arbitrary image into DRAM; it then writes a control register that
switches the UART mux to the NOEL-V and releases the big core's reset. The
NOEL-V begins executing the just-loaded image. This two-core hand-off is the
substrate for everything that follows: every experiment in this paper is an
image streamed to the SERV loader and launched on the NOEL-V. (An intermediate
64-bit core, \emph{Muntjac}, was also explored; a simulation-timeout quirk made
it a poor iteration vehicle, and effort concentrated on the NOEL-V.)

\section{Bringing Up the Core}
\label{sec:core}

\subsection{An X that poisoned the pipeline}
The first milestone was simply getting the NOEL-V to run a boot program and print
``halt.'' It did not: in simulation the core mispredicted branches, speculatively
fetched address~0, and \emph{dropped register writebacks}, so a
\code{lui t0, 0xFF900} became \code{t0=0} and the UART store went to
address~0 instead of \code{0xFF900000}. The cause was neither the fetch unit nor
the store path but an \emph{unconfigured scan interface}: the core's internal AHB
slave record was assembled field by field and never assigned
\reg{testen}/\reg{testrst}/\reg{scanen} (and \reg{endian}). NOEL-V derives its
scan-mux controls from those inputs, so leaving them \code{'U'} propagated X
across flops and RAMs throughout the core. Driving them from the parent bus fixed
it, and the core printed ``halt.'' A general lesson recorded at the time: in
four-state simulation any undriven \code{std\_logic} on a control path becomes X
and spreads far; scan/test inputs are the worst offenders. Diagnose with a VCD
scan for signals stuck at \code{x}/\code{u} after reset.

\subsection{A 64-bit core that was 32-bit}
The next symptom looked like a memory fault: a store to DDR3 ``never reached the
bus''; the bus-interface never asserted its request. Days of tracing the store
path were wasted, because the store was a red herring. An instruction-disassembly
trace showed \emph{32-bit register widths}: the build had compiled
\code{noelv\_cfg\_32.vhd} (\code{NV\_XLEN=32}). An RV32 core running RV64 code
\emph{deadlocked on the first 64-bit word operation} (\code{addw}), several
instructions \emph{before} the store. The \code{cfg} generic (\S\ref{sec:core}
sidebar) selects the microarchitecture family---\code{0x300} is the dual-issue
``GP'' core intended for FPGA---but \emph{not} XLEN, which is the separate global
\code{NV\_XLEN}. Switching to \code{noelv\_cfg\_64.vhd} produced 64-bit registers
and a reachable, working store (AHB monitor: \code{WR 0x40000 hwdata=0xCAFEF00D},
then a matching read-back). Two further self-inflicted wrinkles surfaced in the
tiny hand-written boot assembly: \code{lui t0, 0xFF900} sign-extends to an
invalid high address on RV64 (use \code{li}+\code{slli}), and \code{lw} of a
value with bit~31 set sign-extends and fails a naive compare (use \code{lwu}).

\subsection{A UART that passed simulation and failed on silicon}
With the core running, the on-board serial output was a constant square wave.
The core wrote the correct bytes to the UART data register (confirmed on
debug LEDs capturing the AHB write stream), and the testbench UART decoded
``halt.'' from the transmitter in simulation---yet the GRLIB \code{apbuart\_16550}
transmitter garbled on the real FPGA. Rather than chase a vendored 16550, we
switched to the Gaisler \code{apbuart} (same \reg{uarti}/\reg{uarto} interface, so
no board rewiring), rewrote the boot code for its register map, and set the
scaler for 115200 at the board clock. ``halt.'' then appeared at 115200\,8N1. A
latent cross-domain bug was fixed in passing: a reset synchronizer for the
NOEL-V's reset, released from the SERV clock domain.

\section{The Memory Subsystem}
\label{sec:mem}

\subsection{A byte write that wrote every byte}
Aligned 32- and 64-bit stores worked; single-byte stores (\code{sb}) to DDR3 did
not---and that broke OpenSBI, whose \code{sbi\_memcpy}/\code{sbi\_memset} are
byte-wise. The failure manifested far downstream, as \code{sbi\_domain\_init}
returning \code{SBI\_EINVAL}: a byte-wise region-sort corrupted the root domain's
memory-region table (two swapped entries came back zeroed), and the firmware hung.
A raw bit-bang tracer patched into OpenSBI localized it, and \code{//bin/noelv\_bytetest}
reproduced it in isolation. The bug was in the AHB-to-Wishbone bridge, which
hardwired the Wishbone byte-select strobes to all-ones and never captured the AHB
\reg{hsize}: because the NOEL-V replicates a sub-word datum across all lanes, an
\code{sb} of \code{0xAA} wrote \code{0xAAAAAAAA}. Deriving the strobes from
\reg{hsize} fixed it; a focused testbench (\code{//ip/bridges}) now regression-tests
byte, halfword, word, and burst writes against a lane-honoring model.

\subsection{A 1\,GB memory that looked like 128\,MB}
The board appeared limited to 128\,MB of DRAM. The chip, the UberDDR3 controller,
and the AHB decode were all 1\,GB; the cap was an adapter that dropped byte-address
bits, aliasing the 1\,GB window every 128\,MB. A one-line address-slice fix
restored the full gigabyte, verified by a caches-off march test
(\code{//bin/noelv\_memtest1g}) that writes a distinct value to each 128\,MB bank
and reads all back, and confirmed end to end by a Zircon boot reporting
$\sim$1.0\,GB free after kernel init. This subsystem accreted a small suite of
bare-metal tests---byte, data-section, read/write, atomic (AMO), and
load-reserved/store-conditional probes---each of which paid for itself later.

\section{M-Mode Firmware: OpenSBI}
\label{sec:opensbi}

Once memory was trustworthy, OpenSBI became the focus. Beyond the byte-write heap
corruption above, two core-specific issues blocked the supervisor hand-off.
First, this NOEL-V does not hold \code{lr.d}/\code{sc.d} reservations on uncached
DDR3, so OpenSBI's single use of compare-and-swap (in
\code{atomic\_cmpxchg}) spun forever; since the system is single-hart, a patch
emulates CAS in an interrupt-safe critical section. Second, the core enters
OpenSBI with \reg{mseccfg.MML} set, which the feature detector treats as fatal;
the vendored entry code clears it. With a handful of additional patches (a
drained-putc console so diagnostics are legible on hardware, trusted-FDT parsing
to keep RTL simulation tractable, and trap/init tracers), OpenSBI boots fully:
banner, domain configuration, and a hand-off to a supervisor-mode payload at a
fixed address with the hart ID in \code{a0} and the (expanded) device tree in
\code{a1}.

\section{The Supervisor Boot Ladder}
\label{sec:ladder}

To keep failures localized, the Zircon bring-up was structured as a
\emph{ladder} of increasingly complete bootable images (Table~\ref{tab:ladder}),
each climbed only after the one below was green.

\begin{table}[t]
\caption{The supervisor-mode boot ladder.}
\label{tab:ladder}
\centering
\footnotesize
\begin{tabular}{@{}p{0.05\columnwidth}p{0.55\columnwidth}p{0.26\columnwidth}@{}}
\toprule
\# & What the image proves & Observable \\
\midrule
0 & OpenSBI loads and M-mode runs & OpenSBI banner \\
1 & OpenSBI $\rightarrow$ S-mode hand-off; S-mode executes & ``hello'' payload \\
2 & The Linux boot shim loads, runs in S-mode, parses the DTB & shim diagnostics \\
3 & DTB \code{/chosen} initrd $\rightarrow$ shim finds the ZBI $\rightarrow$ physboot runs & physboot log \\
4 & physboot relocates and starts the kernel $\rightarrow$ userboot & kernel/userboot log \\
\bottomrule
\end{tabular}
\end{table}

Rung~1 revealed a baud mismatch: a payload hardcoded a 50\,MHz scaler, giving
$\sim$92.6\,kbaud on the 40\,MHz board---the banner was transmitted, just
unreadable at a 115200 capture. Rung~2 was subtler. The shim loaded,
self-relocated, and ran, then trapped on an \code{unimp} that
\code{addr2line} placed inside \code{devicetree::Match}: a hardened-libc++
\code{\_\_builtin\_trap} from a \emph{bounds-checked} \code{std::array} iterator,
tripped because our minimal device tree lacked nodes the matchers expected
(a PLIC node, a fuller CPU topology). Completing the device tree cleared it.
Rungs~3--4 fought placement and memory-map issues: physboot placed the data ZBI
at address~0, overwriting the boot stub, until the ZBI load address and the
declared \code{/memory} range were reconciled; caches had to be disabled for the
uncached boot path; a null frame-pointer guard was needed in the backtrace
walker; and the vector extension had to be disabled (the core has no V). The
payload lands at a fixed supervisor address, and the device tree---patched on the
host, since there is no U-Boot---carries the ZBI's location in \code{/chosen}.

The reward was a full boot: OpenSBI $\rightarrow$ shim $\rightarrow$ physboot
$\rightarrow$ the kernel (complete init, memory-pressure subsystem, ``starting
user space'') $\rightarrow$ \code{userboot}, terminating---correctly---at a
deliberate \code{\_\_builtin\_trap} because the engineering ZBI carries no boot
filesystem. This is the same terminus a reference emulator reaches for a
filesystem-less image. An earlier build had one outstanding workaround---with
address-space layout randomization enabled, \code{userboot} took an
instruction-access fault at its randomized high base before running, so the
configuration disabled ASLR---but that no longer reproduces: on the final build
\code{userboot} self-relocates and runs to the same terminus with ASLR enabled,
verified across two independent random placements. The fault had been collateral
of a since-fixed problem (the interrupt-firmware and memory-map changes landed in
between), and ASLR is now left on.

\section{The Interrupt-Driven Console}
\label{sec:console}

Zircon's early output is polled, but its steady-state console is a debug-log
\emph{dumper} thread that performs a \emph{blocking, interrupt-driven}
transmit: it writes a byte and, if the FIFO is full, enables the UART transmit
interrupt and blocks until that interrupt reports space. On our board that thread
never made progress; only the polled fallback
(\code{kernel.debug\_uart\_poll=true}) produced late-boot output. Resolving this
``last mile'' took five independent fixes across four layers, and---as with every
earlier stage---the component everyone suspects, here the platform interrupt
controller (PLIC), was innocent.

\subsection{Two software bugs on the delivery path}
The device tree routed the UART interrupt with the modern
\code{interrupts-extended = <\&plic0 1>} form, but the shim's chosen-node matcher
read only the legacy \code{interrupts} property, resolving the console IRQ to~0
so the kernel never unmasked PLIC source~1. Independently, the kernel PLIC
driver's priority macro computed \code{base + 1 + irq}, writing source~$i$'s
priority into source~$i{+}1$'s slot and leaving source~$i$ at priority~0---below
threshold, hence enabled but never delivered. Both were fixed; neither restored
the console.

\subsection{Exonerating the PLIC}
A JTAG snapshot then showed the PLIC fully and correctly configured for source~1
in the supervisor context, yet \reg{sip.SEIP}${=}0$. Reading the claim register
returned source~1, which we first read as ``the PLIC identifies the interrupt but
fails to drive \reg{seip}''---a hardware bug. An \emph{isolated} simulation of the
\code{grplic\_ahb} block, driven over an AHB bus-functional model exactly as the
kernel drives it, refuted this: given the observed configuration and an asserted
source, the block \emph{does} assert \reg{seip}, and the claim read \emph{clears}
it. The ``\reg{seip}${=}0$ while pending'' reading had been a \emph{measurement
artifact}---the claim read that produced the ``source~1'' datum had itself cleared
the pending interrupt. The PLIC was correct throughout.

\subsection{The core takes no interrupts}
If delivery was fine, why no trap? A bare-metal probe armed the most trivial
interrupt possible---a machine software interrupt---with \reg{mstatus.MIE}${=}1$
and its enable set, verified the pending and enable bits over JTAG, and observed
\emph{no trap}. The same held for timer and external interrupts. Reading the core
RTL resolved it: this NOEL-V implements \emph{Smrnmi}, and the integer unit masks
\emph{all} interrupts while \reg{mnstatus.NMIE}${=}0$---its reset value, which
nothing in the boot chain ever changed. Setting it (\code{csrs 0x744, 8}) let the
bare-metal probe take the full supervisor-external path and claim source~1.
Smrnmi, not the PLIC, was the root cause of ``no interrupts.''

\subsection{A wedge at the supervisor hand-off}
Setting \reg{NMIE} in firmware produced a new failure: the kernel emitted nothing.
With interrupts unmasked, any M-mode interrupt still armed in \reg{mie} fired the
instant OpenSBI executed \code{mret} into supervisor mode---M-interrupts are
always taken in a lower privilege---before the kernel installed its trap vector,
wedging it at entry. Masking \reg{mie} immediately before the \code{mret}
(Listing~\ref{lst:fw}) fixed it; the kernel's own supervisor interrupts are
untouched, and OpenSBI re-arms M-interrupts on demand at its next ecall.

\begin{lstlisting}[caption={The two coupled firmware changes.},label={lst:fw}]
/* fw_base.S, _start_warm, after CSR_MIE is cleared: */
    csrsi   0x744, 8          /* mnstatus.NMIE = 1 */

/* sbi_hart.c, just before the mret to S-mode: */
    csr_write(CSR_MIE, 0);    /* no M-interrupt at S-entry */
\end{lstlisting}

\subsection{A ``stall'' that was really a crawl}
The kernel now booted and \emph{took} timer interrupts (JTAG: \reg{scause} showing
the supervisor timer), yet appeared to hang in threading-level init. The
debug-log ring, dumped over JTAG, told a different story: it was advancing, just
extraordinarily slowly---timestamps put the boot at about nineteen minutes. The
device tree declared \code{riscv,isa = "rv64imac"}, so Zircon could not tell the
core implemented \emph{Sstc} and armed every scheduler-tick timer through an SBI
ecall (a supervisor$\rightarrow$machine$\rightarrow$supervisor round trip) instead
of writing \reg{stimecmp} directly. Advertising the extension
(\code{riscv,isa = "rv64imac\_sstc"}) let the kernel use the direct compare
register; the boot dropped to about seventeen seconds and, with the polled
fallback removed, the \emph{entire} boot streamed over the interrupt-driven
console to the expected terminus, past the init hook at which output had
previously ceased.

\subsection{The five-part fix}
Table~\ref{tab:fixes} collects the changes. They span four layers; no subset
yields a fast, fully interrupt-driven console.

\begin{table}[t]
\caption{The five-part fix for the interrupt console.}
\label{tab:fixes}
\centering
\footnotesize
\begin{tabular}{@{}p{0.05\columnwidth}p{0.55\columnwidth}p{0.26\columnwidth}@{}}
\toprule
\# & Fix & Layer \\
\midrule
1 & Write each source's \emph{own} PLIC priority register & Kernel PLIC driver \\
2 & Resolve the console IRQ from \code{interrupts-extended} & Boot shim \\
3 & Set \reg{mnstatus.NMIE}${=}1$ (unmask on this Smrnmi core) & M-mode firmware \\
4 & Mask \reg{mie} before the S-mode \code{mret} & M-mode firmware \\
5 & Advertise \code{sstc} (fast direct-\reg{stimecmp} timer) & Device tree \\
\bottomrule
\end{tabular}
\end{table}

\section{Debugging Methodology}
\label{sec:method}

Three vantage points, used throughout the arc, made the difference; a
disagreement among them localizes a fault.

\paragraph{Isolated IP simulation}
Simulating one block (the UART transmitter, the AHB bridge, the PLIC) with a bus
functional model observes signals software cannot see and converts ``probably a
hardware bug'' into a proof one way or the other. The PLIC exoneration is the
sharpest example.

\paragraph{Non-perturbing bare-metal probes}
Small OS-free programs arm one interrupt or exercise one memory pattern at a time.
The interrupt probes deliberately read pending state \emph{without} reading the
claim register, precisely because an earlier measurement had been corrupted by
that register's side effect.

\paragraph{A GPL-only JTAG debug module}
The proprietary core's professional debugger refuses to attach, so we drove the
GRLIB JTAG bridge directly from \code{openocd} (GPL) as a plain AHB master:
ground-truth reads of CSRs, of live peripheral registers (PLIC, CLINT), and---most
valuable---of the kernel's in-RAM debug-log ring, which recovered boot progress
when nothing reached the wire. Its limits were charted the hard way (unreliable
single-step; stepping into unloaded RAM wedges the bus), and respected.

\section{Lessons Learned}

\paragraph{Suspect your own non-standard choices first}
An unconfigured scan input, a 32-bit core running 64-bit code, a hardwired byte
strobe, an aliasing address slice, an under-declared device tree---each broke a
component that works for everyone else. The recurring question ``why doesn't
everyone hit this?'' had the same answer every time: because everyone else's
configuration is standard.

\paragraph{Beware self-perturbing measurements}
The ``\reg{seip}${=}0$ while pending'' reading was an artifact of reading the
claim register. On hardware with side-effecting registers, non-invasive probes
are worth the extra care.

\paragraph{Extensions must be enabled \emph{and} advertised}
\emph{Smrnmi} silently masked all interrupts until firmware enabled it;
\emph{Sstc}, present in silicon but absent from \code{riscv,isa}, silently forced
a $100{\times}$ slower timer path. Neither produced a diagnostic; both produced
``it doesn't work'' or ``it hangs.''

\paragraph{Distinguish ``stalled'' from ``slow''}
A log absent from the wire is not necessarily a log that is not being written.
Recovering the in-RAM ring over JTAG and reading its timestamps turned a presumed
deadlock into a measured, then eliminated, performance problem.

\paragraph{Climb a ladder}
Structuring the bring-up as bootable rungs, each adding one layer, meant every new
failure was localized to the layer just added---indispensable when a symptom can
originate in any of five.

\section{Conclusion}

Booting a production microkernel on a RISC-V soft core, end to end on one FPGA,
was a sequence of layer-crossing puzzles whose common thread was that the
suspected component was almost never the culprit. The interrupt controller, the
DRAM, the store path, and the UART core were each exonerated in turn; the real
causes lay in configuration, in a bridge strobe, in firmware that failed to
enable a masking extension, and in a device tree that failed to advertise a timer
extension. A disciplined, multi-vantage methodology---isolating IP in simulation,
probing bare metal without perturbing state, and reading ground truth (including
the kernel's own log) over a GPL JTAG bridge---was what made a low-cost,
fully open toolchain sufficient to find each cause among its several convincing
decoys.

\appendices

\section{Timeline of the Bring-Up}
\label{app:timeline}

Table~\ref{tab:timeline} maps the bring-up to dates, drawn from the project's own
commit history. The RV32/SERV foundation and the DDR3/UART bring-up span the
first year; the NOEL-V, OpenSBI, and Zircon work is concentrated in a dense final
month in which the boot ladder was climbed rung by rung to a full, and then
interrupt-driven, boot.

\begin{table}[ht]
\caption{Timeline of the bring-up.}
\label{tab:timeline}
\centering
\footnotesize
\begin{tabular}{@{}p{0.16\columnwidth}p{0.76\columnwidth}@{}}
\toprule
Period & Milestone \\
\midrule
2024 Q3 & Project begins; SERV RV32 bit-serial core, UART, and DDR3 bring-up in simulation and on the FPGA \\
2024--25 & RV32/SERV boards mature: Intel-HEX serial uploader, DDR3 access and execute-in-place \\
2026-03 & Muntjac RV64 core evaluated as an intermediate \\
2026-06-14 & NOEL-V core introduced \\
2026-06-19 & NOEL-V runs: scan-input X-propagation fixed, RV32$\rightarrow$RV64 config, ``halt.'' on UART (16550$\rightarrow$Gaisler \code{apbuart}) \\
2026-06-21 & Dual-issue NV-HP core builds; OpenSBI \code{fw\_payload} built \\
2026-06-26 & Zircon boot ladder started; shim device-tree bounds-trap root-caused; JTAG debug module (\code{dsuen}) + \code{openocd}/ahbjtag \\
2026-06-27 & OpenSBI boots fully; DDR3 byte-write bridge strobes fixed \\
2026-06-28--30 & Faithful ZBI RTL simulation; a sim-only OpenSBI early-init timing hazard isolated \\
2026-07-05 & Second-simulator (NVC) OpenSBI boot baseline \\
2026-07-07 & DDR3 128\,MB aliasing fixed (adapter address slice); UART 40\,MHz clock \\
2026-07-08 & L1 caches enabled in the boot stub (the Zircon LR/SC fix) \\
2026-07-09 & Full Zircon boot to the expected no-bootfs terminus; PLIC node added \\
2026-07-10 & Full 1\,GB DDR3 verified; entire boot streamed on the polled serial console \\
2026-07-13 & Interrupt-driven console solved (Smrnmi \reg{NMIE} + \reg{mie}-quiesce + \code{sstc}) \\
\bottomrule
\end{tabular}
\end{table}

\end{document}
